

\section{El último teorema de Fermat}

\[
7x^2 + x - 5 = 0
\]
\[
x^2 + y^2 = z^2
\]
\[
x^3 + y^3 = z^3
\]
\[
x^4 + y^4 = z^4
\]

\begin{nota}{Descenso infinito}
Fermat demostró que si $(x,y,z)$ es una solución en enteros positivos de
$x^4+y^4=z^4$, entonces se puede construir otra solución con números
más chicos, y así obtener una cadena infinita descendente de soluciones
cada vez menores, lo cual es imposible en $\mathbb{N}$ (método de
\textbf{descenso infinito} de Fermat). Esto prueba que la ecuación no
tiene soluciones no triviales.
\end{nota}

De manera más general, la ecuación de Fermat (o de Diofanto) es
\[
x^n + y^n = z^n, \qquad n \geq 3,
\]
y el último teorema de Fermat afirma que no tiene soluciones enteras
positivas.

\section{Propiedades de la suma y el producto en $\mathbb{Z}$}
\begin{teoria}
El conjunto $\mathbb{Z}$ viene provisto de dos operaciones, $+$ y $*$.

\textbf{Suma:}
\begin{listapropiedades}
    \item[Asociativa] $(x+y)+z = x+(y+z)$. Por ejemplo,
    $(1+4)+7 = 1+(4+7) = 12$.

    \item[Conmutativa] $x+y = y+x$. La suma de varios términos se puede
    reordenar libremente, por ejemplo $1+2+3+4 = 4+3+2+1$.

    \item[Elemento neutro] Tiene a $0$ como neutro: $x+0=x$ para todo
    $x \in \mathbb{Z}$.

    \item[Elemento opuesto] Todo entero tiene un opuesto (inverso
    aditivo): para cada $x \in \mathbb{Z}$ existe $-x \in \mathbb{Z}$
    tal que $x+(-x)=0$.
\end{listapropiedades}

\textbf{Multiplicación:}
\begin{listapropiedades}
    \item[Asociativa] $x*(y*z) = (x*y)*z$.

    \item[Conmutativa] $x*y = y*x$.

    \item[Elemento neutro] Tiene a $1$ como neutro: $x*1 = x$.

    \item[Sin inverso general] En general, un entero $x \neq 0$
    \emph{no} tiene inverso multiplicativo en $\mathbb{Z}$ (salvo
    $x=\pm 1$): no siempre existe $y \in \mathbb{Z}$ tal que $x*y=1$.
\end{listapropiedades}

\textbf{Propiedad distributiva:}
\[
x*(y+z) = x*y + x*z.
\]
De ahí se deduce, por ejemplo,
\[
(x+y)*z = x*z + y*z.
\]

Con estas propiedades, $(\mathbb{Z},+,*)$ es un \textbf{anillo
conmutativo con unidad}.
\end{teoria}

\begin{definicion}
Un anillo conmutativo con unidad $(R,+,*)$ es un \textbf{cuerpo} si todo
elemento distinto de $0$ tiene inverso multiplicativo: para cada
$x \in R\setminus\{0\}$ existe $y \in R$ tal que $x*y=1$.
\end{definicion}

\begin{nota}{$\mathbb{Z}$ no es un cuerpo}
$\mathbb{Z}$ no es un cuerpo (por ejemplo $2$ no tiene inverso
multiplicativo entero), pero sí es un \textbf{dominio de integridad}:
no tiene divisores de cero, es decir
\[
xy = 0 \;\Longrightarrow\; x=0 \text{ o } y=0.
\]
\end{nota}

\begin{ejemplo}
$21$ es divisible por $7$, ya que $21 = 7\cdot 3$.
\end{ejemplo}

\section{Divisibilidad}

\begin{definicion}
Sean $a,b \in \mathbb{Z}$. Decimos que $a$ \textbf{divide} a $b$ si
existe un entero $c$ tal que $b=ac$, y en ese caso escribimos $a\mid b$.
Si $a$ no divide a $b$ escribimos $a \nmid b$.
\end{definicion}

\begin{proposicion}
Consideremos la relación
\[
R = \{ (a,b) \in \mathbb{Z} \times \mathbb{Z} : a \mid b \}.
\]


$R$ es reflexiva y transitiva, pero no es antisimétrica (y por lo tanto
no es una relación de orden en $\mathbb{Z}$).
\end{proposicion}

\begin{demostracion}
\textbf{Reflexiva.} Sea $a \in \mathbb{Z}$. Como $a=a\cdot 1$, se tiene
$a\mid a$, así que $(a,a) \in R$.

\textbf{No simétrica.} $R$ no es simétrica: $(3,9) \in R$ (pues
$3\mid 9$) pero $(9,3) \notin R$ (pues $9 \nmid 3$).

\textbf{Transitiva.} Supongamos que $(a,b),(b,c) \in R$, es decir
$a\mid b$ y $b \mid c$. Entonces existen enteros $x,y$ tales que $b=ax$
y $c=by$. Luego
\[
c = by = (ax)y = a(xy),
\]
así que $a \mid c$, es decir $(a,c) \in R$. Por lo tanto $R$ es
transitiva.

\textbf{No antisimétrica.} Supongamos que $(a,b),(b,a) \in R$, es decir
que hay enteros $x,y$ tales que $b=ax$ y $a=by$. Entonces
\[
a = by = (ax)y = a(xy),
\]
de donde
\[
0 = a - a(xy) = a(1-xy),
\]
así que $a=0$ o $xy=1$.

\begin{itemize}
    \item Si $a=0$, entonces $b=ax=0$, luego $a=b$.
    \item Si $a\neq 0$, entonces $xy=1$ con $x,y \in \mathbb{Z}$, así
    que $x$ es invertible en $\mathbb{Z}$, es decir $x=1$ o $x=-1$.
    Si $x=1$, entonces $b=ax=a$. Si $x=-1$, entonces $b=ax=-a$.
\end{itemize}

En conclusión, $(a,b),(b,a)\in R$ implica $a=b$ o $a=-b$, pero no
necesariamente $a=b$: por ejemplo $a=1$, $b=-1$ satisface $1\mid -1$ y
$-1\mid 1$, pero $a\neq b$. Por lo tanto $R$ \textbf{no} es
antisimétrica en $\mathbb{Z}$.
\end{demostracion}

\begin{proposicion}
La relación
\[
R' = \{ (a,b) \in \mathbb{N} \times \mathbb{N} : a \mid b \}
\]
sí es una relación de orden (parcial) en $\mathbb{N}$.
\end{proposicion}

\begin{demostracion}
En $\mathbb{N}$ el caso $a=-b$ no puede ocurrir (salvo $a=b=0$), así que
de $(a,b),(b,a) \in R'$ se deduce $a=b$, y $R'$ resulta antisimétrica,
además de reflexiva y transitiva (por el mismo argumento de la
demostración anterior).
\end{demostracion}

\section{Congruencia módulo $m$}

\begin{definicion}
Fijemos un entero $m$. Decimos que un entero $a$ es \textbf{congruente
módulo $m$} con otro entero $b$ si $m \mid a-b$, y en ese caso
escribimos
\[
a \equiv b \pmod m.
\]
\end{definicion}

\begin{teoria}
Sea
\[
R = \{ (a,b) \in \mathbb{Z}\times\mathbb{Z} : a \equiv b \pmod m \}.
\]
$R$ es una relación en $\mathbb{Z}$.

La relación de congruencia módulo $m$ es una relación de equivalencia
en $\mathbb{Z}$.
\end{teoria}

\begin{demostracion}
\textbf{Reflexiva.} Sea $a \in \mathbb{Z}$. Como $m \mid 0 = a-a$, se
tiene $a \equiv a \pmod m$, así que $(a,a) \in R$.

\textbf{Simétrica.} Supongamos $(a,b) \in R$, es decir $a\equiv b
\pmod m$. Entonces existe $x \in \mathbb{Z}$ tal que $a-b=mx$. Luego
\[
b-a = -(a-b) = -mx = m(-x),
\]
así que $m \mid b-a$, es decir $b \equiv a \pmod m$ y $(b,a)\in R$.

\textbf{Transitiva.} Supongamos $(a,b),(b,c) \in R$, es decir
$a\equiv b \pmod m$ y $b \equiv c \pmod m$. Entonces existen enteros
$x,y$ tales que
\[
a-b=mx, \qquad b-c=my.
\]
Sumando ambas igualdades,
\[
(a-b)+(b-c) = mx+my,
\]
es decir
\[
a-c = m(x+y).
\]
Como $x+y \in \mathbb{Z}$, se tiene $m \mid a-c$, así que
$a \equiv c \pmod m$ y $(a,c) \in R$.

Por lo tanto $R$ es reflexiva, simétrica y transitiva, es decir, es una
relación de equivalencia en $\mathbb{Z}$.
\end{demostracion}

sea R = {(a,b) in Z X Z : a== B mod m}
R es una relacion en Z
si a in Z , entonces m | 0 = a - a , asi que a== a y (a,a) in R . R es reflexiva
$
(a,b) \rightarrow (b,a) in R 
a== b , m|a-b 
a - b = mx 
b-a = -mx 
= m(-x)
b==a 
m|b-a 
b-a = m (-x)
si (a,b) in R , entonces a== b asi que existe x in R talque a-b =mx entonces m(-x) = -mx = -(a-b) = b-a , asi que m|b-a , b==a , asi que (b,a) in R . R es simetrica
(a,b),(b,c) e R -> (a,c) e R 
a==b b==c a==c
a-b= mx b-c=my a-c=mc
a = b+mx
c = b-my
mz = a-c = (b+mx) -(b-my)
mz = mx + my
si que (a,b) y (b,c) esta en R, de manera que ahy entero x e y talque  a-b = mx y b-c=my 
entonces m(x+y) = mx = my = (a-b) + (b-c)
= a-b + b -c 
= a-c
asi que m|a-c
y a==c , (a,c) in R
== es una relacion de equivalencia en Z$